\setupcolors[state=start]
\setupcolor[hex]
\definecolor[darkblue][h=003366]
\definecolor[darkgray][h=e6e6e6]
\definecolor[redline][h=d9534f]
\definecolor[blueline][h=5bc0de]
\definecolor[greenline][h=5cb85c]

\setuphead[section][style=\bfa, color=darkblue, before=\blank, after=\blank]
\setuphead[subsection][style=\bf, color=black, before=\blank, after=\blank]
\setuppapersize[A4][A4]
\setupbodyfont[11pt]
\setuplayout[width=middle, height=middle, topspace=2cm, bottomspace=2cm]
\mainlanguage[es]

\starttext
\title{\bf Comparación de OPP-Split\crlf con OPP-HC y OPP-ACO}
\blank[big]

\startsection[title={Resumen}]
\startitemize
\item {\bf Las soluciones de las dos heurísticas son válidas}: las %%NUM:recomputed_checked%%
verificadas contra el modelo monolítico cumplen todas las restricciones del modelo exacto, y el objetivo
recalculado coincide con el reportado en todas.
\item {\bf Coinciden con el exacto en el estado}: donde split encuentra solución, las heurísticas
también; donde split prueba que no hay, tampoco la encuentran.
\item {\bf ACO-4 es casi exacto}: alcanza el óptimo en las %%NUM:optimal_functional_opp_aco%%
instancias funcionales comparables y queda a %%NUM:max_gap_stress_opp_aco%%~\% como máximo en
stress.
\item {\bf OPP-HC es el más rápido}: en stress es %%NUM:stress_speedup_opp_hc%% veces más rápido
que split en la mediana, con una brecha de %%NUM:max_gap_stress_opp_hc%%~\% como máximo.
\item {\bf Escalan donde split no}: en los %%NUM:stress_noref%% escenarios de stress que split no
termina en 40 minutos, las dos heurísticas dan una solución válida.
\stopitemize
\stopsection

\startsection[title={Objetivo y alcance}]
Este documento compara el modelo exacto OPP-Split (formulación en dos etapas, resuelta con
SCIP) con las dos metaheurísticas que resuelven el mismo problema: Hill Climbing (OPP-HC, en las
tablas OPP-HC) y Ant Colony Optimization (OPP-ACO, ACO-4). Las tres resuelven ruteo, asignación y
secuenciación con velocidades heterogéneas de pickers y cadena de frío.

Las preguntas son cinco: si las soluciones de las heurísticas son válidas para el modelo exacto,
cuánto se alejan del óptimo, cuánto más rápido llegan, cómo escalan con el tamaño de la instancia
y cuánta memoria usan.

Los datos son los CSV consolidados:
\startitemize
\item \type{results/comparability/results_split_vs_heuristics.csv} — una fila por instancia y
heurística, con el estado, el objetivo y el tiempo de split y de la heurística, la brecha y el
cociente de tiempos. Lo escribe \type{python -m experiments.comparison.comparability} a partir de los CSV de
resultados de los tres modelos.
\item \type{results/opp_hc/stress/}, \type{results/opp_aco/stress/} y
\type{results/opp_split/stress/} — tamaño de cada instancia, memoria, nivel de servicio y
convergencia.
\item \type{results/compliance/opp_hc/} y \type{results/compliance/opp_aco/} — la verificación
de cada solución de las heurísticas contra las restricciones del modelo exacto.
\stopitemize
\stopsection

\startsection[title={Los tres métodos}]
\blank[medium]
\bTABLE[split=yes, align=middle, frame=on, width=fit]
\setupTABLE[row][1][background=color, backgroundcolor=darkgray, style=bold]
\setupTABLE[column][3][width=6.5cm, align=flushleft]
\bTR \bTH Método \eTH \bTH Tipo \eTH \bTH Configuración en stress \eTH \bTH Tope \eTH \eTR
\bTR \bTD OPP-Split \eTD \bTD Exacto (ZIMPL + SCIP) \eTD \bTD Etapa 1: un MIP de ruteo por lote, en 6 procesos. Etapa 2: asignación y secuenciación. \eTD \bTD 40 min por test \eTD \eTR
\bTR \bTD OPP-HC \eTD \bTD Búsqueda local \eTD \bTD Hasta %%NUM:hc_max_it%% iteraciones; corta tras %%NUM:hc_max_imp%% sin mejora. \eTD \bTD 30 s de reloj, mirado entre iteraciones \eTD \eTR
\bTR \bTD ACO-4 \eTD \bTD Colonia de hormigas \eTD \bTD %%NUM:aco_ants%% hormigas, %%NUM:aco_max_it%% iteraciones, $\alpha = %%NUM:aco_alpha%%$, $\beta = %%NUM:aco_beta%%$, %%NUM:aco_route_attempts%% intentos de ruta por lote. \eTD \bTD %%NUM:aco_max_time_min%% min de reloj, mirado entre iteraciones \eTD \eTR
\eTABLE
\blank[medium]

Los topes de las heurísticas no interrumpen una iteración en curso: se comparan con el reloj
al terminar cada una. En el 01 (100×100) eso las lleva a %%NUM:hc_01_min%% y
%%NUM:aco_01_min%% minutos: OPP-HC supera sus 30~s ya en la construcción de la solución inicial, y
cada iteración de ACO-4 dura más de un cuarto de hora.

Las heurísticas corren una sola vez por instancia, con semilla 1. Los tiempos son de reloj,
medidos en la misma notebook, e incluyen la preparación de cada test (en split, escribir los
\type{.dat} y compilar ZIMPL).
\stopsection

\startsection[title={Qué se compara}]
Se comparan solo instancias {\em equivalentes}: los tres modelos resuelven exactamente los mismos
datos. En la suite funcional son las instancias \type{test_instance__*} de las heurísticas, que
reproducen cada escenario que resuelve un test funcional de split. En stress son los 19 escenarios
de la suite, que los tres modelos construyen con el mismo generador y los mismos argumentos. Las
reglas están en \type{docs/specs/comparabilidad_split_heuristicas.md}; en resumen:
\startitemize
\item {\bf Estado y objetivo}: split encontró el óptimo y no usa espera voluntaria.
\item {\bf Solo estado}: split declara la instancia infactible (se verifica que la heurística
tampoco encuentre solución), o split usa {\em espera voluntaria}. En el Split un lote puede figurar
como terminado más tarde de lo que termina físicamente, para reducir la penalización por
adelanto; las heurísticas reportan tiempos reales y no pueden hacerlo, así que ahí el objetivo de
split es menor por diseño del modelo y no por un plan mejor.
\item {\bf No comparable}: split no termina, o el test no llega al solver.
\stopitemize

%%TABLE:coverage%%

Las funcionales que quedan fuera:

%%TABLE:functional_out%%

Los %%NUM:stress_out%% escenarios de stress no comparables son aquellos en los que split no
termina en 40 minutos: no hay referencia de objetivo. Se usan para ver escala (sección
\in[sec:escala]) y para comparar las heurísticas entre sí (sección \in[sec:noref]).
\stopsection

\startsection[title={Validez de las soluciones heurísticas}]
Cada solución de las heurísticas se verificó contra OPP-Mono: se fijan los valores de la
solución sobre las variables del modelo monolítico y SCIP decide si son factibles, sin copiar ni
reescribir ninguna restricción (\type{verification/run_compliance.py}). También se recalcula el
objetivo y se compara con el que reporta la heurística.

%%TABLE:compliance%%

Ninguna solución viola una restricción del modelo exacto, y el objetivo recalculado coincide con el reportado
en %%NUM:recomputed_matched%% de las %%NUM:recomputed_checked%% verificadas. Las instancias {\em sin
solución} son las infactibles: la heurística no devuelve solución y no hay nada que verificar. Por
lo tanto, las brechas que siguen comparan soluciones válidas del mismo problema, con el objetivo
bien calculado.
\stopsection

\startsection[title={Estado}]
Instancias en las que la heurística coincide con split en el estado: encuentra solución donde split
la encuentra y no la encuentra donde split prueba que la instancia es infactible.

%%TABLE:status%%

Las dos heurísticas coinciden en todas las instancias comparables en estado. En las
%%NUM:functional_infeasible%% infactibles no devuelven solución, como corresponde; en las demás
siempre encuentran una. Hay que leerlo con cuidado: una heurística no prueba infactibilidad, solo no
encuentra solución. Que coincida con split en las infactibles muestra que no inventa soluciones
inválidas, y que en las factibles nunca se queda sin respuesta.
\stopsection

\startsection[title={Suite funcional}]

\startsubsection[title={Calidad de la solución}]
La brecha es $100 \cdot (z_{\rm heur} - z_{\rm split}) / z_{\rm split}$, sobre las instancias
comparables en objetivo. {\em Óptimo} cuenta las instancias con brecha nula. Ninguna heurística
queda por debajo del óptimo de split, como corresponde a soluciones válidas de un problema de
minimización.

%%TABLE:gap%%

\startitemize
\item {\bf ACO-4} encuentra el óptimo en las %%NUM:optimal_functional_opp_aco%% instancias
funcionales.
\item {\bf OPP-HC} encuentra el óptimo en %%NUM:optimal_functional_opp_hc%% de las
%%NUM:functional_both%% funcionales.
\stopitemize

La brecha en stress está en la sección \in[sec:brecha_stress].

Las %%NUM:missed%% instancias funcionales en las que alguna heurística no llega al óptimo son
todas de grilla, y en todas es OPP-HC. Una búsqueda local se queda ahí en un óptimo local: la grilla
tiene muchas rutas y asignaciones de costo parecido, y los movimientos de a uno no alcanzan para
salir.

%%TABLE:missed%%
\stopsubsection

\startsubsection[title={Instancias con espera voluntaria}]
En %%NUM:wait_count%% instancias funcionales la solución de split demora a propósito algún lote:
lo reporta terminado más tarde de lo que termina, para acercarlo a la fecha de entrega y reducir
la penalización por adelanto. Las heurísticas reportan el tiempo en que el lote realmente termina y
no pueden hacerlo, así que ahí el objetivo de split es menor por diseño del modelo, no por un plan
mejor. Por eso estas instancias se comparan solo en estado, y la diferencia no es una brecha:

%%TABLE:wait%%

ACO-4 queda por encima de split en %%NUM:wait_aco_above%% de las %%NUM:wait_count%%. La
especificación de comparabilidad registra que, en estas instancias, tomar las rutas más cortas y
agregarles la espera óptima reproduce exactamente el objetivo de split: la diferencia es la espera.
\stopsubsection

\startsubsection[title={Tiempos}]
Sobre las %%NUM:timed%% instancias comparables en objetivo cuyo tiempo de split se conoce: el CSV
funcional de split guarda una fila por test, así que el tiempo de una instancia se toma solo
cuando el test resuelve una sola instancia.

%%TABLE:functional_times%%

\placefigure[force][fig:funcional]{Tiempo total de cada instancia funcional: split contra cada heurística (escala logarítmica en los dos ejes). Debajo de la diagonal, la heurística fue más rápida.}
{%%FIG:functional_scatter%%}

Las instancias funcionales son chicas y los tiempos de split, del orden de 0,3~s, están
dominados por el costo fijo de escribir y compilar los modelos. OPP-HC es
%%NUM:functional_speedup_opp_hc%% veces más rápido en la mediana y es más lento que split en
%%NUM:functional_slower_opp_hc%% de las %%NUM:timed%% instancias. ACO-4 es
%%NUM:functional_speedup_opp_aco%% veces más rápido en la mediana, pero más lento que split en
%%NUM:functional_slower_opp_aco%% de las %%NUM:timed%% (los puntos por encima de la diagonal,
cerca de 1~s): corre sus iteraciones completas aunque la instancia sea trivial.
\stopsubsection
\stopsection

\startsection[title={Suite de stress}]

\startsubsection[title={Las instancias}]
El perfil realista (26 órdenes por lote, dos bandas de velocidad) en grillas de 7×7 a 16×16 y
cuadrillas de hasta 24 lotes, más tres escenarios de grilla (01 a 03). Los tres
modelos resuelven la misma instancia en cada escenario.

%%TABLE:sizes%%

En los de grilla el peso por lote completado $W^c$ es 9 y todas las órdenes llegan tarde por
construcción; en el perfil realista es 27, mayor que las órdenes por lote, así que la espera
voluntaria nunca conviene y los objetivos son comparables.
\stopsubsection

\startsubsection[title={Tiempos}]
%%TABLE:stress%%

\placefigure[force][fig:stress]{Tiempo total por escenario de stress (escala logarítmica). Las barras sin relleno son los timeouts de split.}
{%%FIG:stress_times%%}

En los %%NUM:stress_both%% escenarios que split resuelve, OPP-HC es %%NUM:stress_speedup_opp_hc%%
veces más rápido que split en la mediana y ACO-4, %%NUM:stress_speedup_opp_aco%% veces. OPP-HC
resuelve todos menos el 01 en menos de 5~s; ACO-4 tarda entre 10~s y 10 minutos, del mismo orden que
split en los escenarios chicos y bastante menos en los grandes.

El tiempo de split en el 10 es el de la corrida del 26 de septiembre (86~s). La corrida más
reciente, con el mismo objetivo, marcó 412~s porque la notebook se suspendió durante el test.
\stopsubsection

\startsubsection[title={Escalabilidad}, reference=sec:escala]
El mismo tiempo, ahora contra el tamaño de la instancia (lotes × ubicaciones de la grilla).

\placefigure[force][fig:escala]{Tiempo total contra el tamaño de la instancia (escala logarítmica en los dos ejes). Los círculos abiertos son los escenarios que split no termina en 40 minutos.}
{%%FIG:stress_scale%%}

Split no tiene un tamaño límite único: termina la grilla de 15×15 con 6 lotes pero no la de
16×16, y con 8 lotes o más no termina ni en 7×7, donde el producto lotes × ubicaciones es chico.
Son dos paredes distintas: la Etapa 1 crece con la grilla y la Etapa 2 con la cantidad de lotes.
Las heurísticas no tienen ninguna de las dos: OPP-HC se mantiene en segundos en todo el rango salvo
el 100×100, y ACO-4 en minutos.
\stopsubsection

\startsubsection[title={Calidad de la solución}, reference=sec:brecha_stress]
%%TABLE:gap_stress%%

\placefigure[force][fig:brecha]{Brecha contra el óptimo de split en cada escenario que split resuelve. Un 0 indica que la heurística encontró el óptimo.}
{%%FIG:stress_gap%%}

ACO-4 queda por debajo de OPP-HC en %%NUM:gap_aco_below_hc%% de los %%NUM:stress_solved%%
escenarios. Las dos brechas son chicas: la mayor, %%NUM:max_gap_stress_opp_hc%%~\% de OPP-HC, es
en el escenario con cadena de frío (11).
\stopsubsection

\startsubsection[title={Calidad contra tiempo}]
Cada punto es un modelo en un escenario que split resuelve. Abajo a la izquierda está lo
deseable: rápido y cerca del óptimo.

\placefigure[force][fig:tradeoff]{Brecha contra tiempo total en los escenarios de stress que split resuelve (tiempo en escala logarítmica).}
{%%FIG:tradeoff%%}

Los tres modelos ocupan zonas distintas: split prueba el óptimo (brecha 0) pagando el mayor
tiempo; ACO-4 queda prácticamente en el óptimo con algo menos de tiempo; OPP-HC resigna unas
centésimas de punto de brecha a cambio de uno o dos órdenes de magnitud de tiempo. Ninguno domina
a los otros en las dos dimensiones a la vez.
\stopsubsection

\startsubsection[title={Escenarios sin referencia}, reference=sec:noref]
En los %%NUM:stress_noref%% escenarios que split no termina, las dos heurísticas encuentran
solución. No hay óptimo con qué medirlas, pero sí se pueden comparar entre sí:

%%TABLE:noref%%

ACO-4 obtiene el mejor objetivo en %%NUM:noref_aco_better%% de los %%NUM:stress_noref%%, con
diferencias de menos de 1~\%. La excepción es el 01 (30 lotes en 100×100): ACO-4 alcanza su
tope de %%NUM:aco_max_time_min%% minutos después de solo %%NUM:aco_01_iterations%% iteraciones, y
la notebook se suspendió durante esa corrida, así que su objetivo ahí vale como solución factible,
no como medida de calidad.
\stopsubsection
\stopsection

\startsection[title={Nivel de servicio}]
Porcentaje de órdenes entregadas a tiempo en cada escenario de stress, según la solución de cada
modelo.

%%TABLE:sla%%

En los de grilla (01 a 03) ninguna orden llega a tiempo en ningún modelo, por construcción del
escenario. En el perfil realista los tres entregan todo a tiempo, salvo en el 11 (cadena de frío),
donde el límite de frío obliga a demorar algunas órdenes: ahí ACO-4 iguala el nivel de servicio del
óptimo y OPP-HC queda levemente por debajo.
\stopsection

\startsection[title={Memoria}]
Memoria pico (RAM) del proceso de cada test, en los escenarios de stress.

\placefigure[force][fig:memoria]{Memoria pico por escenario de stress (escala logarítmica). Split no tiene dato en los timeouts.}
{%%FIG:memory%%}

Las heurísticas usan unos %%NUM:mem_median_opp_hc%%~MB (OPP-HC) y %%NUM:mem_median_opp_aco%%~MB
(ACO-4) en la mediana; en el 100×100 (01) suben a %%NUM:mem_01_opp_hc%% y
%%NUM:mem_01_opp_aco%%~MB por el tamaño del grafo. ACO-4 tiene picos de hasta
%%NUM:mem_aco_spike_max%%~MB en algunos escenarios (%%NUM:mem_aco_spikes%%), que no siguen al
tamaño de la instancia. Split usa entre %%NUM:mem_split_min%% y %%NUM:mem_split_max%%~MB, pero
esa cifra subestima su consumo: mide solo el proceso del test, sin los seis procesos de la Etapa 1
ni el de ZIMPL.
\stopsection

\startsection[title={Convergencia de las heurísticas}]
Iteraciones que corrió cada heurística en stress y en cuál encontró la mejor solución (0 es la
solución inicial).

%%TABLE:convergence%%

\startitemize
\item {\bf OPP-HC} corta tras %%NUM:hc_max_imp%% iteraciones sin mejora, así que corre entre 60 y
160 iteraciones. En %%NUM:hc_best_at_start%% escenarios la solución inicial ya era la mejor que
encontró. Solo en el 03 y el 17 llega al máximo de %%NUM:hc_max_it%% iteraciones. En el 01 no llega
a iterar: la construcción inicial del 100×100 ya supera su tope de 30~s.
\item {\bf ACO-4} corre sus %%NUM:aco_max_it%% iteraciones. La mejor solución aparece en cualquier
punto, de la iteración 1 a la 54, y en la mitad de los escenarios después de la 30: más iteraciones
podrían mejorarla algo. Solo en el %%NUM:aco_timed_out%% corta por tiempo.
\stopitemize
\stopsection

\startsection[title={Limitaciones}]
\startitemize
\item Una corrida por instancia, con semilla 1: no se mide la variabilidad de las heurísticas
entre semillas.
\item Los tiempos son de reloj en una notebook; si se suspende durante una corrida, la duración
incluye el tiempo suspendido. Afecta sobre todo al 01 de ACO-4.
\item La referencia de objetivo existe solo donde split termina: %%NUM:stress_both%% de los
%%NUM:stress_total%% escenarios de stress. En los otros, la calidad de las heurísticas solo se
compara entre ellas.
\item La memoria de split no incluye sus procesos hijos.
\stopitemize
\stopsection

\startsection[title={Conclusiones}]
\startitemize[n]
\item {\bf Validez:} todas las soluciones de las dos heurísticas cumplen las restricciones del modelo exacto,
verificadas contra el modelo monolítico, y reportan bien su objetivo.
\item {\bf Estado:} las dos heurísticas coinciden con split en las %%NUM:status_comparable%%
instancias funcionales comparables en estado y en los %%NUM:stress_both%% escenarios de stress
que split resuelve, incluidas las infactibles.
\item {\bf Calidad:} ACO-4 alcanza el óptimo en todas las funcionales y queda a
%%NUM:max_gap_stress_opp_aco%%~\% como máximo en stress. OPP-HC falla el óptimo en
%%NUM:missed%% instancias de grilla (hasta %%NUM:max_gap_functional_opp_hc%%~\%) y queda a
%%NUM:max_gap_stress_opp_hc%%~\% como máximo en stress.
\item {\bf Tiempo:} en stress, OPP-HC es el más rápido (%%NUM:stress_speedup_opp_hc%% veces más
rápido que split en la mediana); ACO-4 lo es %%NUM:stress_speedup_opp_aco%% veces, a cambio de
una brecha menor.
\item {\bf Escala:} donde split no termina en 40 minutos (%%NUM:stress_noref%% escenarios), las
dos heurísticas dan una solución válida: OPP-HC en segundos y ACO-4 en minutos, salvo el 01
(100×100), que les lleva 17 y 35 minutos.
\item {\bf Recursos:} las heurísticas usan decenas de MB salvo en el grafo de 100×100; ACO-4
tiene picos puntuales de más de 1~GB.
\stopitemize

{\bf Recomendación.} Hasta 6 lotes en grillas de hasta 15×15, split da el óptimo probado en
minutos y es la referencia. Por encima de eso, ACO-4 es la opción cuando importa la calidad: queda
prácticamente en el óptimo y da siempre una solución válida. OPP-HC conviene cuando manda el tiempo
de respuesta, por ejemplo para replanificar en línea: resigna centésimas de punto de calidad a
cambio de responder en segundos.
\stopsection

\stoptext
